\documentclass[11pt, a4paper]{article}

\usepackage[T1]{fontenc}
\usepackage[utf8]{inputenc}
\usepackage{tgpagella}
\usepackage[spanish, es-noshorthands]{babel}
\usepackage{amsmath, amssymb}
\usepackage[most]{tcolorbox}
\usepackage[american]{circuitikz}
\usepackage[margin=2.5cm]{geometry}
\usepackage{fancyhdr}

\pagestyle{fancy}
\fancyhf{}
\setlength{\headheight}{16pt}
\lhead{\textbf{UCB Sede Tarija} | Ing. Mecatrónica}
\rhead{IMT-221 | Diagnóstico S06-2}
\renewcommand{\headrulewidth}{0.4pt}
\definecolor{ucbblue}{RGB}{0, 51, 102}

\begin{document}

\begin{tcolorbox}[colback=ucbblue!5, colframe=ucbblue, title=\textbf{Diagnóstico Breve de Prerrequisitos DC}]
    Tiempo estimado: 5-7 minutos. Resuelve individualmente para habilitar el análisis AC.
\end{tcolorbox}

\section*{Ejercicio 1: Evaluación del Punto Q}
Considera el siguiente circuito de polarización fija. 
Datos: $V_{CC}=10\,\text{V}$, $\beta=100$, $V_{BE} \approx 0.7\,\text{V}$ (en conducción).

\begin{center}
\begin{circuitikz}[scale=1, transform shape, thick]
    \draw (0,0) node[npn](Q){};
    \draw (Q.E) node[ground]{};
    \draw (Q.C) to[R, l_={$R_C = 2.0\,\text{k}\Omega$}] (0,3);
    \draw (Q.B) -- (-2.5, 0) to[R, l={$R_B = 470\,\text{k}\Omega$}] (-2.5, 3);
    \draw (-2.5,3) -- (0,3) -- (0,3.5) node[above]{$+10\,\text{V}$};
\end{circuitikz}
\end{center}

\textbf{Calcula y responde:}
\begin{enumerate}
    \item Calcula la corriente de base $I_B$. \vspace{1.5cm}
    \item Calcula la corriente candidata de colector $I_C$. \vspace{1.5cm}
    \item Calcula el voltaje $V_{CE}$. \vspace{1.5cm}
    \item Verifica matemáticamente la consistencia de la región (¿está en Activa Directa?). \vspace{1.5cm}
    \item Escribe el punto de operación exacto $Q = (I_{CQ}, V_{CEQ})$. \vspace{1cm}
    \item En una sola oración técnica: ¿Qué representa físicamente el punto $Q$? \vspace{1.5cm}
    \item ¿Por qué el punto $Q$ debe conocerse \textit{antes} de realizar el análisis de pequeña señal?
\end{enumerate}

\end{document}
